% Options for packages loaded elsewhere
% Options for packages loaded elsewhere
\PassOptionsToPackage{unicode}{hyperref}
\PassOptionsToPackage{hyphens}{url}
\PassOptionsToPackage{dvipsnames,svgnames,x11names}{xcolor}
%
\documentclass[
  letterpaper,
  DIV=11,
  numbers=noendperiod]{scrartcl}
\usepackage{xcolor}
\usepackage[top=1in, bottom=0.75in, left=1in, right=1in,
footskip=0.5in]{geometry}
\usepackage{amsmath,amssymb}
\setcounter{secnumdepth}{-\maxdimen} % remove section numbering
\usepackage{iftex}
\ifPDFTeX
  \usepackage[T1]{fontenc}
  \usepackage[utf8]{inputenc}
  \usepackage{textcomp} % provide euro and other symbols
\else % if luatex or xetex
  \usepackage{unicode-math} % this also loads fontspec
  \defaultfontfeatures{Scale=MatchLowercase}
  \defaultfontfeatures[\rmfamily]{Ligatures=TeX,Scale=1}
\fi
\usepackage{lmodern}
\ifPDFTeX\else
  % xetex/luatex font selection
\fi
% Use upquote if available, for straight quotes in verbatim environments
\IfFileExists{upquote.sty}{\usepackage{upquote}}{}
\IfFileExists{microtype.sty}{% use microtype if available
  \usepackage[]{microtype}
  \UseMicrotypeSet[protrusion]{basicmath} % disable protrusion for tt fonts
}{}
\makeatletter
\@ifundefined{KOMAClassName}{% if non-KOMA class
  \IfFileExists{parskip.sty}{%
    \usepackage{parskip}
  }{% else
    \setlength{\parindent}{0pt}
    \setlength{\parskip}{6pt plus 2pt minus 1pt}}
}{% if KOMA class
  \KOMAoptions{parskip=half}}
\makeatother
% Make \paragraph and \subparagraph free-standing
\makeatletter
\ifx\paragraph\undefined\else
  \let\oldparagraph\paragraph
  \renewcommand{\paragraph}{
    \@ifstar
      \xxxParagraphStar
      \xxxParagraphNoStar
  }
  \newcommand{\xxxParagraphStar}[1]{\oldparagraph*{#1}\mbox{}}
  \newcommand{\xxxParagraphNoStar}[1]{\oldparagraph{#1}\mbox{}}
\fi
\ifx\subparagraph\undefined\else
  \let\oldsubparagraph\subparagraph
  \renewcommand{\subparagraph}{
    \@ifstar
      \xxxSubParagraphStar
      \xxxSubParagraphNoStar
  }
  \newcommand{\xxxSubParagraphStar}[1]{\oldsubparagraph*{#1}\mbox{}}
  \newcommand{\xxxSubParagraphNoStar}[1]{\oldsubparagraph{#1}\mbox{}}
\fi
\makeatother


\usepackage{longtable,booktabs,array}
\usepackage{calc} % for calculating minipage widths
% Correct order of tables after \paragraph or \subparagraph
\usepackage{etoolbox}
\makeatletter
\patchcmd\longtable{\par}{\if@noskipsec\mbox{}\fi\par}{}{}
\makeatother
% Allow footnotes in longtable head/foot
\IfFileExists{footnotehyper.sty}{\usepackage{footnotehyper}}{\usepackage{footnote}}
\makesavenoteenv{longtable}
\usepackage{graphicx}
\makeatletter
\newsavebox\pandoc@box
\newcommand*\pandocbounded[1]{% scales image to fit in text height/width
  \sbox\pandoc@box{#1}%
  \Gscale@div\@tempa{\textheight}{\dimexpr\ht\pandoc@box+\dp\pandoc@box\relax}%
  \Gscale@div\@tempb{\linewidth}{\wd\pandoc@box}%
  \ifdim\@tempb\p@<\@tempa\p@\let\@tempa\@tempb\fi% select the smaller of both
  \ifdim\@tempa\p@<\p@\scalebox{\@tempa}{\usebox\pandoc@box}%
  \else\usebox{\pandoc@box}%
  \fi%
}
% Set default figure placement to htbp
\def\fps@figure{htbp}
\makeatother





\setlength{\emergencystretch}{3em} % prevent overfull lines

\providecommand{\tightlist}{%
  \setlength{\itemsep}{0pt}\setlength{\parskip}{0pt}}



 


\KOMAoption{captions}{tableheading}
\makeatletter
\@ifpackageloaded{tcolorbox}{}{\usepackage[skins,breakable]{tcolorbox}}
\@ifpackageloaded{fontawesome5}{}{\usepackage{fontawesome5}}
\definecolor{quarto-callout-color}{HTML}{909090}
\definecolor{quarto-callout-note-color}{HTML}{0758E5}
\definecolor{quarto-callout-important-color}{HTML}{CC1914}
\definecolor{quarto-callout-warning-color}{HTML}{EB9113}
\definecolor{quarto-callout-tip-color}{HTML}{00A047}
\definecolor{quarto-callout-caution-color}{HTML}{FC5300}
\definecolor{quarto-callout-color-frame}{HTML}{acacac}
\definecolor{quarto-callout-note-color-frame}{HTML}{4582ec}
\definecolor{quarto-callout-important-color-frame}{HTML}{d9534f}
\definecolor{quarto-callout-warning-color-frame}{HTML}{f0ad4e}
\definecolor{quarto-callout-tip-color-frame}{HTML}{02b875}
\definecolor{quarto-callout-caution-color-frame}{HTML}{fd7e14}
\makeatother
\makeatletter
\@ifpackageloaded{caption}{}{\usepackage{caption}}
\AtBeginDocument{%
\ifdefined\contentsname
  \renewcommand*\contentsname{Table of contents}
\else
  \newcommand\contentsname{Table of contents}
\fi
\ifdefined\listfigurename
  \renewcommand*\listfigurename{List of Figures}
\else
  \newcommand\listfigurename{List of Figures}
\fi
\ifdefined\listtablename
  \renewcommand*\listtablename{List of Tables}
\else
  \newcommand\listtablename{List of Tables}
\fi
\ifdefined\figurename
  \renewcommand*\figurename{Figure}
\else
  \newcommand\figurename{Figure}
\fi
\ifdefined\tablename
  \renewcommand*\tablename{Table}
\else
  \newcommand\tablename{Table}
\fi
}
\@ifpackageloaded{float}{}{\usepackage{float}}
\floatstyle{ruled}
\@ifundefined{c@chapter}{\newfloat{codelisting}{h}{lop}}{\newfloat{codelisting}{h}{lop}[chapter]}
\floatname{codelisting}{Listing}
\newcommand*\listoflistings{\listof{codelisting}{List of Listings}}
\makeatother
\makeatletter
\makeatother
\makeatletter
\@ifpackageloaded{caption}{}{\usepackage{caption}}
\@ifpackageloaded{subcaption}{}{\usepackage{subcaption}}
\makeatother
\usepackage{bookmark}
\IfFileExists{xurl.sty}{\usepackage{xurl}}{} % add URL line breaks if available
\urlstyle{same}
\hypersetup{
  pdftitle={Multi-Sector Models and Marketing Margins},
  pdfauthor={Lauren Chenarides},
  colorlinks=true,
  linkcolor={blue},
  filecolor={Maroon},
  citecolor={Blue},
  urlcolor={Blue},
  pdfcreator={LaTeX via pandoc}}


\title{Multi-Sector Models and Marketing Margins}
\author{}
\date{}
\begin{document}
\maketitle


\subsection{Elasticities and
Forecasting}\label{elasticities-and-forecasting}

Elasticities are useful tools in applied economics. They tell us:

\begin{itemize}
\tightlist
\item
  whether a tax on a specific good will raise large or small revenues,\\
\item
  which goods are the greatest competitors, and\\
\item
  how prices and quantities may change in response to shocks.
\end{itemize}

For example:

\begin{itemize}
\tightlist
\item
  How does the supply of pork decreases due to tougher regulations on
  manure treatment and disposal?\\
\item
  How much will the demand for pork increase due to an increase in the
  price of beef?
\end{itemize}

We do not need to know the \emph{exact} supply and demand curves.
Estimates of elasticities are usually sufficient.

By incorporating elasticities into an \textbf{equilibrium displacement
model (EDM)}, we can calculate price and quantity changes resulting from
outside shocks.

\subsection{Endogenous and Exogenous
Variables}\label{endogenous-and-exogenous-variables}

Equilibrium displacement models are used to assess how \emph{endogenous
variables} respond to changes in \emph{exogenous variables}.

In an EDM, all variables are grouped as either \textbf{endogenous} or
\textbf{exogenous}.

\begin{itemize}
\tightlist
\item
  \textbf{Endogenous variables}: Determined \emph{within} the model
  (price and quantity).\\
\item
  \textbf{Exogenous variables}: Determined \emph{outside} the model and
  taken as given.

  \begin{itemize}
  \tightlist
  \item
    Examples: input prices, technology, consumer income, population,
    tastes, and prices of related goods.
  \end{itemize}
\end{itemize}

Shocks to exogenous variables shift supply and/or demand curves. The
market then finds a new equilibrium where supply equals demand.

\subsection{Formalizing the Model}\label{formalizing-the-model}

\subsubsection{Demand Curve}\label{demand-curve}

\paragraph{Step 1. Start with the elasticity
definition}\label{step-1.-start-with-the-elasticity-definition}

For demand, the own-price elasticity is defined as:

\[
E_D = \frac{\% \Delta Q_D}{\% \Delta P}
\]

This formula says: ``a 1\% change in price leads to \(E_D\) percent
change in quantity demanded.''\\
Remember that \(E_D\) is typically \textbf{negative} because price and
quantity demanded move in opposite directions.

\paragraph{Step 2. Rearrange to express quantity
changes}\label{step-2.-rearrange-to-express-quantity-changes}

Multiply both sides by \(\% \Delta P\) to solve for the percent change
in quantity demanded:

\[
\% \Delta Q_D = E_D \cdot (\% \Delta P)
\]

This equation tells us how much quantity demanded changes in response to
a given price change, holding everything else constant.

\paragraph{Step 3. Add in demand
shocks}\label{step-3.-add-in-demand-shocks}

Markets can also shift for reasons other than price (income, tastes,
population, etc.).\\
We capture these \textbf{exogenous demand shifts} with a new term,
\(S_D\):

\[
\% \Delta Q_D = E_D (\% \Delta P) + S_D
\]

\begin{itemize}
\tightlist
\item
  \(\% \Delta Q_D\): percent change in quantity demanded\\
\item
  \(\% \Delta P\): percent change in price\\
\item
  \(E_D\): own-price elasticity of demand (negative)\\
\item
  \(S_D\): demand shock (e.g., if incomes rise and demand increases by
  7\%, then \(S_D = 7\))
\end{itemize}

\textbf{Takeaway:} This step-by-step process shows how we move from the
basic elasticity definition to an equation that includes both
\textbf{price effects} and \textbf{shifts in demand.}

\subsubsection{Supply Curve}\label{supply-curve}

\paragraph{Step 1. Start with the elasticity
definition}\label{step-1.-start-with-the-elasticity-definition-1}

For supply, the own-price elasticity is defined as:

\[
E_S = \frac{\% \Delta Q_S}{\% \Delta P}
\]

This formula says: ``a 1\% change in price leads to \(E_S\) percent
change in quantity supplied.''\\
Unlike demand, \(E_S\) is typically \textbf{positive} because price and
quantity supplied move in the same direction.

\paragraph{Step 2. Rearrange to express quantity
changes}\label{step-2.-rearrange-to-express-quantity-changes-1}

Multiply both sides by \(\% \Delta P\) to solve for the percent change
in quantity supplied:

\[
\% \Delta Q_S = E_S \cdot (\% \Delta P)
\]

This equation captures how much quantity supplied changes in response to
a price change, holding everything else constant.

\paragraph{Step 3. Add in supply
shocks}\label{step-3.-add-in-supply-shocks}

Supply can also shift due to exogenous factors (e.g., technology, input
costs, regulations).\\
We capture these \textbf{exogenous supply shifts} with a new term,
\(S_S\):

\[
\% \Delta Q_S = E_S (\% \Delta P) + S_S
\]

\begin{itemize}
\tightlist
\item
  \(\% \Delta Q_S\): percent change in quantity supplied\\
\item
  \(\% \Delta P\): percent change in price\\
\item
  \(E_S\): own-price elasticity of supply (positive)\\
\item
  \(S_S\): supply shock (e.g., if new technology lowers costs and supply
  increases 5\%, then \(S_S = 5\))
\end{itemize}

\subsection{Solving for the Equilibrium
Price}\label{solving-for-the-equilibrium-price}

Now we combine supply and demand.

\begin{enumerate}
\def\labelenumi{\arabic{enumi}.}
\item
  In equilibrium, quantity supplied equals quantity demanded:

  \[
  \% \Delta Q_S = \% \Delta Q_D
  \]
\item
  Substitute in the supply and demand equations:

  \[
  E_S(\% \Delta P) + S_S = E_D(\% \Delta P) + S_D
  \]
\item
  Rearrange terms:

  \[
  (E_S - E_D)(\% \Delta P) = S_D - S_S
  \]
\item
  Solve for the endogenous variable, price:

  \[
  \% \Delta P = \frac{S_D - S_S}{E_S - E_D}
  \]
\end{enumerate}

\subsection{Intuition}\label{intuition}

\begin{itemize}
\tightlist
\item
  The denominator \((E_S - E_D)\) is positive since \(E_S > 0\) and
  \(E_D < 0\).\\
\item
  If \(S_D > 0\) (a positive demand shock, e.g., higher incomes) and
  \(S_S = 0\), then \(\% \Delta P > 0\) → price rises.\\
\item
  If \(S_S > 0\) (a positive supply shock, e.g., lower production costs)
  and \(S_D = 0\), then \(\% \Delta P < 0\) → price falls.
\end{itemize}

\textbf{Takeaway:} With both curves expressed in elasticity form, we can
solve for how price responds to shocks. Once \(\% \Delta P\) is known,
we substitute back into either curve to find the equilibrium change in
quantity.

\subsection{Solving for the Equilibrium
Quantity}\label{solving-for-the-equilibrium-quantity}

Once we know the equilibrium price change \(\% \Delta P\), we can solve
for the change in quantity.

\begin{enumerate}
\def\labelenumi{\arabic{enumi}.}
\tightlist
\item
  Substitute \(\% \Delta P\) into the supply curve
\end{enumerate}

Recall the supply curve:

\[
\% \Delta Q_S = E_S (\% \Delta P) + S_S
\]

Plug in the solved value for \(\% \Delta P\):

\[
\% \Delta Q_S = E_S \left( \frac{S_D - S_S}{E_S - E_D} \right) + S_S
\]

\begin{enumerate}
\def\labelenumi{\arabic{enumi}.}
\setcounter{enumi}{1}
\tightlist
\item
  Or substitute into the demand curve
\end{enumerate}

Recall the demand curve:

\[
\% \Delta Q_D = E_D (\% \Delta P) + S_D
\]

Substitute the same \(\% \Delta P\):

\[
\% \Delta Q_D = E_D \left( \frac{S_D - S_S}{E_S - E_D} \right) + S_D
\]

\begin{enumerate}
\def\labelenumi{\arabic{enumi}.}
\setcounter{enumi}{2}
\tightlist
\item
  Confirm consistency
\end{enumerate}

In equilibrium, both give the same answer:

\[
\% \Delta Q_S = \% \Delta Q_D = \% \Delta Q
\]

If they don't match, a calculation error has been made.

\subsection{Intuition}\label{intuition-1}

\begin{itemize}
\tightlist
\item
  A positive demand shock (\(S_D > 0\)) raises price and increases
  equilibrium quantity.\\
\item
  A positive supply shock (\(S_S > 0\)) lowers price but still increases
  equilibrium quantity.\\
\item
  The exact size of the change depends on how responsive supply and
  demand are (their elasticities).
\end{itemize}

\textbf{Takeaway:} First solve for \(\% \Delta P\) using equilibrium,
then plug it into either curve to get \(\% \Delta Q\). Both equations
must agree at the new equilibrium.

\subsection{Quick Recap}\label{quick-recap}

Once \(\% \Delta P\) is found, substitute back into either curve:

\begin{itemize}
\tightlist
\item
  Supply:\\
  \[
  \% \Delta Q_S = E_S(\% \Delta P) + S_S
  \]
\item
  Demand:\\
  \[
  \% \Delta Q_D = E_D(\% \Delta P) + S_D
  \]
\end{itemize}

Both give the same \(\% \Delta Q\) in equilibrium.

\subsection{General Steps for an EDM}\label{general-steps-for-an-edm}

\begin{enumerate}
\def\labelenumi{\arabic{enumi}.}
\tightlist
\item
  \textbf{Identify shocks:} Determine \(S_S\) and \(S_D\). At least one
  must be nonzero.\\
\item
  \textbf{Specify supply curve:}\\
  \[
  \% \Delta Q_S = E_S (\% \Delta P) + S_S
  \]
\item
  \textbf{Specify demand curve:}\\
  \[
  \% \Delta Q_D = E_D (\% \Delta P) + S_D
  \]
\item
  \textbf{Solve for price:} Set \(\% \Delta Q_S = \% \Delta Q_D\) and
  solve for \(\% \Delta P\).\\
\item
  \textbf{Solve for quantity:} Substitute \(\% \Delta P\) into either
  equation to calculate \(\% \Delta Q\).
\end{enumerate}

\newpage

\subsection{Example: Impacts of Manure Regulations in the Pork
Market}\label{example-impacts-of-manure-regulations-in-the-pork-market}

Pork producers are facing tougher regulations on how they manage and
treat hog manure. These regulations raise production costs, which shifts
the supply curve \textbf{leftward}, raising price and lowering quantity.
We quantify these changes using elasticities and an equilibrium
displacement model (EDM).

\begin{tcolorbox}[enhanced jigsaw, bottomtitle=1mm, colbacktitle=quarto-callout-note-color!10!white, opacitybacktitle=0.6, colback=white, colframe=quarto-callout-note-color-frame, rightrule=.15mm, leftrule=.75mm, left=2mm, breakable, title=\textcolor{quarto-callout-note-color}{\faInfo}\hspace{0.5em}{Note}, toptitle=1mm, titlerule=0mm, arc=.35mm, bottomrule=.15mm, toprule=.15mm, opacityback=0, coltitle=black]

This example uses \textbf{long-run elasticities} (\(E_S = 2.15\),
\(E_D = -1.96\)).\\
To examine the short-run impacts, simply replace these with short-run
elasticities.

\end{tcolorbox}

\subsubsection{Step 1. Express the cost increase as a percent of the
current
price}\label{step-1.-express-the-cost-increase-as-a-percent-of-the-current-price}

Suppose the current hog price is \$50/cwt, and the new regulations
increase production costs by \$1/cwt.

\[
P = \$50 \ \text{per cwt} 
\] \[
\Delta C = \$1 \ \text{per cwt}
\]

Calculate the cost increase as a percentage of the current price:

\[
\frac{\Delta C}{P} = \frac{\$1}{\$50} = 0.02 = 2\%
\]

\textbf{Interpretation:} Paying \$1 more per cwt is like receiving \$1
less per cwt in revenue. From the producer's perspective, this is an
\textbf{effective price decrease} of 2\%:

\[
\% \Delta P_{\text{effective}} = -2\%.
\]

\subsubsection{Step 2. Convert the effective price change into a supply
shock}\label{step-2.-convert-the-effective-price-change-into-a-supply-shock}

The \textbf{supply elasticity} tells us how responsive producers are to
a change in price:

\[
E_S = \frac{\% \Delta Q_S}{\% \Delta P}.
\]

Rearrange to solve for the percent change in quantity supplied:

\[
\% \Delta Q_S = E_S \times (\% \Delta P).
\] This equation says: if the effective price producers receive changes
by some percent, then quantity supplied will change by elasticity × that
percent. In our case, the regulation is equivalent to a \textbf{2\% fall
in effective price}. It is also important to know that the (long-run)
supply elasticity \(E_S = 2.15\)

\newpage

Substitute the effective price change from Step 1 (\(-2\%\)) and the
long-run supply elasticity \(E_S = 2.15\):

\[
S_S = E_S \times (\% \Delta P_{\text{effective}})
\] \[
S_S = 2.15 \times (-2)
\] \[
S_S = -4.3\%
\]

\textbf{Interpretation:} A 2\% decrease in effective price leads to a
4.3\% decrease in supply.\\
This is the \textbf{exogenous supply shock}: \(S_S = -4.3\%\).\\
Because nothing changed on the demand side, \(S_D = 0\).

\subsubsection{Step 3. Write the EDM
equations}\label{step-3.-write-the-edm-equations}

Supply:

\[
\% \Delta Q_S = E_S (\% \Delta P) + S_S = 2.15(\% \Delta P) - 4.3
\]

Demand:

\[
\% \Delta Q_D = E_D (\% \Delta P) + S_D = -1.96(\% \Delta P) + 0
\]

\subsubsection{Step 4. Impose Equilibrium (Set supply equal to
demand)}\label{step-4.-impose-equilibrium-set-supply-equal-to-demand}

In equilibrium, \(\% \Delta Q_S = \% \Delta Q_D\):

\[
2.15(\% \Delta P) - 4.3 = -1.96(\% \Delta P)
\]

Rearrange:

\[
(2.15 + 1.96)(\% \Delta P) = 4.3
\]

Solve for price:

\[
\% \Delta P = \frac{4.3}{2.15+1.96} = 1.05 \%.
\]

\textbf{Result:} Price rises by about \textbf{1.05\%}.

\subsubsection{Step 5. Solve for
Quantity}\label{step-5.-solve-for-quantity}

Plug back into supply or demand:

Using supply:

\[
\% \Delta Q_S = 2.15(1.05) - 4.3
\] \[
\% \Delta Q_S \approx -2.05\%
\]

Check with demand:

\[
\% \Delta Q_D = -1.96(1.05) \approx -2.05\%
\] \textbf{Result:} Quantity falls by about \textbf{2.05\%}.

\subsection{Pork Market: Supply Shift from Manure
Regulations}\label{pork-market-supply-shift-from-manure-regulations}

\begin{figure}[h!]
\centering
\begin{tikzpicture}[>=Latex, scale=1.05]

% Axes
\draw[->] (0,0) -- (9,0) node[below right] {$Q$ (pork)};
\draw[->] (0,0) -- (0,6) node[left] {$P$ (\$/cwt)};

% Demand line D
\coordinate (Dleft)  at (0.8,5.3);
\coordinate (Dright) at (8.2,0.7);
\draw[thick, blue] (Dleft) -- (Dright) node[pos=0.9, below, xshift=6pt] {$D$};

% Initial Supply S0
\coordinate (S0left)  at (0.8,0.9);
\coordinate (S0right) at (8.2,4.9);
\draw[thick, green!60!black] (S0left) -- (S0right) node[pos=0.85, above, xshift=8pt] {$S_0$};

% New Supply S1
\coordinate (S1left)  at ($(S0left)+( -1.1, 0.6)$);
\coordinate (S1right) at ($(S0right)+(-1.1, 0.6)$);
\draw[thick, green!60!black, dashed] (S1left) -- (S1right) node[pos=0.85, above, xshift=8pt] {$S_1$};

% Intersections
\path[name path=D]    (Dleft) -- (Dright);
\path[name path=S0]   (S0left) -- (S0right);
\path[name path=S1]   (S1left) -- (S1right);

\path[name intersections={of=D and S0, by=E0}];
\path[name intersections={of=D and S1, by=E1}];

% Guides
\draw[densely dotted] (E0) -- ($(E0 |- 0,0)$) node[below] {$Q_0$};
\draw[densely dotted] (E0) -- ($(0,0 |- E0)$) node[left]  {$P_0$};

\draw[densely dotted] (E1) -- ($(E1 |- 0,0)$) node[below] {$Q_1$};
\draw[densely dotted] (E1) -- ($(0,0 |- E1)$) node[left]  {$P_1$};

% Equilibrium points
\fill (E0) circle (1.6pt) node[below right] {$E_0$};
\fill (E1) circle (1.6pt) node[above left]  {$E_1$};

\end{tikzpicture}
\caption{Manure regulations raise costs $\Rightarrow$ supply shifts left ($S_0 \to S_1$): $P$ rises, $Q$ falls.}
\end{figure}

\begin{tcolorbox}[enhanced jigsaw, bottomtitle=1mm, colbacktitle=quarto-callout-note-color!10!white, opacitybacktitle=0.6, colback=white, colframe=quarto-callout-note-color-frame, rightrule=.15mm, leftrule=.75mm, left=2mm, breakable, title=\textcolor{quarto-callout-note-color}{\faInfo}\hspace{0.5em}{Interpretation}, toptitle=1mm, titlerule=0mm, arc=.35mm, bottomrule=.15mm, toprule=.15mm, opacityback=0, coltitle=black]

\begin{itemize}
\item
  Although regulations raised production costs by 2\%, the market price
  rose by only about 1\%. In other words, a 2\% cost increase (relative
  to price) cannot be fully passed on to consumers.
\item
  The remainder of the adjustment occurred through reduced output: pork
  production fell by about 2\%.
\item
  This reflects how elasticities distribute the burden of shocks between
  \textbf{price changes} and \textbf{quantity changes}.
\end{itemize}

\end{tcolorbox}




\end{document}
